\documentclass[10pt,twoside]{article}
\usepackage[utf8]{inputenc}
\usepackage{amsmath}
\usepackage{amsfonts}
\usepackage{amssymb}
\usepackage[spanish,es-noshorthands]{babel}
\usepackage[T1]{fontenc}
\usepackage{lmodern}
\usepackage{graphicx,hyperref}
\usepackage{tikz,pgf}
\usepackage{multicol}
\usepackage{subfig}
\usepackage[papersize={6.5in,8.5in},width=5.5in,height=7in]{geometry}
\usepackage{fancyhdr}
\pagestyle{fancy}
\fancyhead[LE]{\includegraphics[height=12pt]{Images/logo-colegio.png} Geometría $9^{\circ}$}
\fancyhead[RE]{}
\fancyhead[RO]{matematicas.german@gmail.com}
\fancyhead[LO]{}

\author{Germ\'an Avenda\~no Ram\'irez,~\thanks{Lic. Mat. U.D., M.Sc. U.N.}}
\title{\begin{minipage}{.2\textwidth}
\includegraphics[height=1.75cm]{Images/logo-colegio.png}\end{minipage}
\begin{minipage}{.55\textwidth}
\begin{center}
Taller 01\\
Proporcionalidad\\
Geometr\'{i}a $9^{\circ}$
\end{center}
\end{minipage}\hfill
\begin{minipage}{.2\textwidth}
\includegraphics[height=1.75cm]{Images/logo-sed.png} 
\end{minipage}}
\date{}
\begin{document}
\maketitle
Nombre: \hrulefill Curso: \underline{\hspace*{44pt}} Fecha: \underline{\hspace*{2.5cm}}
\paragraph*{Meta de aprendizaje}: El estudiante aplica las propiedades de las proporciones para hallar incógnitas
\subsection*{Entre uno y otro}
\begin{enumerate}
\item Empleando una expresi\'{o}n matem\'{a}tica escribir simb\'{o}licamente las siguientes expresiones
\begin{enumerate}
\item Tengo 4 tortas para tres personas
\item Se fue ocho d\'{i}as y solo trabaj\'{o} cuatro.
\item Camilo da tres pasos en tres segundos
\item Realice la notaci\'{o}n de las estrategias empleadas a partir de la Ruleta Heur\'{i}stica.
\end{enumerate}
\end{enumerate}
\section*{Raz\'{o}n matem\'{a}tica}
La razón matemática es una expresión que se encarga de relacionar dos cantidades sin tener en cuenta el tipo. Se representa mediante un cociente indicado, $\dfrac{a}{b}$ \'{o} $a:b$. La lectura es $a$ es a $b$, también se puede emplear $a$ de $b$.
\subsection*{Partes de la raz\'{o}n}
 \[\dfrac{a}{b}=\dfrac{Antecedente}{Consecuente}\]
 La razón suele expresarse en fracción reducida (simplificada), sus partes son el antecedente $a$ y  el consecuente $b$.
 \subsubsection*{Algunos ejemplos}
 \begin{enumerate}
 \item 5 es a 7, es la razón entre días de asistencia al colegio y días de la semana.
 \item 1 es 3, razón entre cada uno de los colores de la bandera de Colombia y el total de colores de la bandera. $\dfrac{1}{3}$ o $1:3$
 \end{enumerate}
 \subsection*{Escribiendo razones}
 \begin{enumerate}
 \item Escriba las siguientes razones, empleando las cuatro formas de expresar una razón, y en el caso que sea posible hallar la fracción reducida.
 \begin{enumerate}
 \item 15 galletas con trocitos de menta en una bolsa con 34 galletas.
 \item 16 perros pastor alemán de 24 perros.
 \item 25 conjuntos residenciales de ladrillo de cada 45.
 \item 10 tambores de 75 instrumentos.
 \item 32 vacas de 72 mamíferos.
 \item Realice la notación de las estrategias empleadas a partir de la Ruleta Heurística.
 \end{enumerate}
 \end{enumerate}
 \subsection*{De razones a tasas}
 \begin{enumerate}
 \item Expresar en forma de razón las siguientes expresiones.
 \begin{enumerate}
 \item 8 kilómetros recorre en una  hora.
 \item 30 vueltas da un disco en un minuto.
 \item c) 45 kilogramos en una botella de 3 litros.
 \item 35 personas por metro cuadrado.
 \item Establecer las diferencias y semejanzas con las razones del numeral dos.
 \item 32 vacas de 72 mamíferos.
 \item Realice la notación de las estrategias empleadas a partir de la Ruleta Heurística.
 \end{enumerate}

 \end{enumerate}
\end{document}
